\chapter{Background}
\label{sec:background}
% TODO: briefly explain / mention the types of compressors here. Maybe explain
% in Compression Methods below.

\section{History of Compression}
\label{sec:compression-history}
First, we will go over a brief summary and history of compression methods. More
detail about said methods will be given in~\cref{sec:compression-methods}. In
the beginning, the words were uncompressed and it was good.\cite{salomon_data_2004}
\subsection{Entropy Encoders}
\label{subsec:entropy-compressor-history}
In 1952, \citeauthor{huffman_method_1952} published
\citetitle{huffman_method_1952} which introduced a coding method which became
known as Huffman coding (HC). Huffman coding improved upon Shannon-Fano coding.
More specifics of the algorithm are discussed in~\cref{sec:huffman-coding}.

As early as 1999, arithmetic coding (AC) was used in the best compression
methods for text and several state-of-the-art image compression
algorithms.\cite[p.110]{mackay_information_2003} In 2014,
\citeauthor{duda_asymmetric_2013}
% TODO: short title or full title?
published \citetitle{duda_asymmetric_2013},\footnote{Should this be the
  full-title? Right now, it's the ``short title'' stored in Zotero.} introducing
a new approach to entropy coding: asymmetric numeral systems (ANS). This method
``provides about 50\% faster decoding than HC for 256 size alphabet, with
compression rate similar to provided by AC.''\cite{duda_asymmetric_2013}.

\subsection{Dictionary-Based Compressors}
\label{subsec:dict-compressor-history}
In 1977 and 1978, LZ77 and LZ78, respectively, were released by
\citeauthor{ziv_universal_1977} in their papers
\citetitle{ziv_universal_1977}\cite{ziv_universal_1977} and
\citetitle{ziv_compression_1978}\cite{ziv_compression_1978}. These algorithms
(and their descendents, known as Lempel-Ziv (LZ) compressors) are
dictionary-based compressors rather than entropy (en)coders.

\section{Compression Methods}
\label{sec:compression-methods}

\subsection{Huffman Coding}
\label{sec:huffman-coding}

